\clearpage
\appendix

\section{Finite Algorithm}\label{app:algorithm}

\noindent\fbox{\begin{minipage}{0.94\textwidth}
\textbf{Algorithm 1: Source-scored structural initial design}
\begin{enumerate}[label=\arabic*.]
\item Generate the declared 64-profile outcome-free library on 128 regular
midpoints using seed 20260808.
\item For every source task and profile, compute the replicated objective mean
and the Gaussian chance-margin statistic in \eqref{eq:source-margin}.
\item Convert each statistic to deterministic within-task average-tie
percentile ranks; average the ranks uniformly across source tasks.
\item Compute exact Voronoi-cell cosine coefficients $c_0,\ldots,c_8$, divide
$c_k$ by $1+.25k$, append all nine diagonal squares, and standardize the 18
coordinates over the library.
\item Append source safety and objective ranks.  Select the first center by
\eqref{eq:first-center}; select nine further centers by
\eqref{eq:farthest-first}.
\item Linearly interpolate each selected profile to the declared target grid,
clip to bounds, and deterministically round.  Freeze these ten policies before
the first target outcome.
\item Run any declared target backend.  Freeze a distinct ordered shortlist
before verification.
\item Test each shortlist candidate on its independent all-success stream.
Deploy the first certificate; abstain if none certifies.
\end{enumerate}
\end{minipage}}

The stable tie order is safety rank, objective rank, profile identifier for the
first center and finite-library index for farthest-first.  Robust source
feasibility is diagnostic only.  There is no endpoint replacement, hidden
scalarization, target-outcome repair, or target-oracle branch in a non-oracle
arm.

\section{Additional Statistical Results}\label{app:additional-results}

\subsection{Verifier thresholds}

Table \ref{tab:verifier-power-table} gives exact power calculations.  At
$n=80$, the Clopper--Pearson lower-bound rule equals the all-success rule.  At
$n=160$, its preregistered sensitivity threshold permits two failures.  The
primary results retain the all-success protocol.

\begin{table}[htbp]
\centering
\small
\caption{Exact candidate-wise certificate power at budgets 80 and 160.}
\label{tab:verifier-power-table}
\begin{tabular}{rrrrrr}
\toprule
$n$ & $p$ & All-success power & CP threshold & Allowed failures & CP power \\
\midrule
80 & 0.950 & 0.017 & 80 & 0 & 0.017 \\
80 & 0.975 & 0.132 & 80 & 0 & 0.132 \\
80 & 0.990 & 0.448 & 80 & 0 & 0.448 \\
80 & 0.995 & 0.670 & 80 & 0 & 0.670 \\
80 & 1.000 & 1.000 & 80 & 0 & 1.000 \\
160 & 0.950 & 0.000 & 158 & 2 & 0.012 \\
160 & 0.975 & 0.017 & 158 & 2 & 0.234 \\
160 & 0.990 & 0.200 & 158 & 2 & 0.784 \\
160 & 0.995 & 0.448 & 158 & 2 & 0.953 \\
160 & 1.000 & 1.000 & 158 & 2 & 1.000 \\
\bottomrule
\end{tabular}

\end{table}

\subsection{Full sensitivity matrix}

The full one-factor sensitivity matrix is reported in Table
\ref{tab:sensitivity-full}.  Entries are ordered by the tested levels in the
second column.  Effective-rank levels are explicit overrides; all other rows
include the frozen baseline level.  The nonmonotone responses, the zero
certificate rate at $\alpha=.01$, and the immaterial frequency penalty are
retained rather than optimized away.

\begin{table}[htbp]
\centering
\scriptsize
\caption{One-factor sensitivity of the source-scored design at $d=1000$.
Each rate aggregates eight regimes and 20 independent tasks per regime.}
\label{tab:sensitivity-full}
\resizebox{0.98\textwidth}{!}{\begin{tabular}{llll}
\toprule
Factor & Tested levels & Feasible (\%) & Certified (\%) \\
\midrule
Effective rank & 4/8/16/32 & 89.4/85.0/84.4/83.8 & 55.6/36.3/25.6/29.4 \\
Chance level $\alpha$ & .01/.05/.1 & 95.0/95.0/95.0 & 0.0/46.2/29.4 \\
Calibrated safe mass & .03/.08/.18 & 69.4/95.0/99.4 & 22.5/46.2/70.0 \\
Initial-design size $n_0$ & 5/10/20 & 79.4/95.0/96.9 & 42.5/46.2/46.9 \\
Source tasks & 1/2/4 & 85.0/95.0/95.6 & 33.8/46.2/45.0 \\
Profiles per source task & 32/64/128 & 83.1/95.0/73.8 & 38.1/46.2/38.8 \\
Source replications & 1/3/5 & 92.5/95.0/92.5 & 45.0/46.2/45.6 \\
Retained frequency $K$ & 4/8/16 & 81.9/95.0/95.6 & 43.1/46.2/48.1 \\
Frequency penalty $\kappa$ & 0/.25/1 & 95.0/95.0/95.0 & 46.2/46.2/46.2 \\
Safety-rank weight & .25/1/4 & 95.0/95.0/95.0 & 46.9/46.2/46.2 \\
Objective-rank weight & .25/1/4 & 95.0/95.0/95.0 & 46.9/46.2/46.2 \\
First-center safety weight & .25/.5/.75 & 62.5/95.0/95.0 & 32.5/46.2/45.0 \\
\bottomrule
\end{tabular}
}
\end{table}

Recorded wall time is descriptive rather than a resource-normalized endpoint.
On the mixed CPU cluster, the primary ten-point initial-design arms average
1.2--1.7 seconds per synthetic task, whereas 394-call source-free controls
average about 59 seconds and the 394-call functional-SCBO arm has a median of
about 677 seconds.  The synthetic simulator is analytic and cheap, so these
figures measure implementation overhead rather than the time of a physical
simulator; simulation-call accounting remains the portable comparison.

\subsection{Native end-to-end transfer}

\begin{table}[htbp]
\centering
\small
\caption{Native transfer pipelines at $N=20$ with the same 384-call source
archive and independent verifier.  Each method retains its own initialization
and online adaptation.}
\label{tab:native-transfer}
\resizebox{0.88\textwidth}{!}{\begin{tabular}{lrrrrr}
\toprule
Native method & FactorShock & Inventory & Queue & Total & False \\
\midrule
Safe F-PACOH & 0/20 & 0/20 & 0/20 & 0/60 & 0 \\
RGPE & 14/20 & 16/20 & 0/20 & 30/60 & 0 \\
Stacked GP & 0/20 & 13/20 & 0/20 & 13/60 & 0 \\
MTGP & 1/20 & 0/20 & 0/20 & 1/60 & 0 \\
FSBO & 6/20 & 2/20 & 0/20 & 8/60 & 0 \\
HyperBO & 0/20 & 0/20 & 0/20 & 0/60 & 0 \\
MetaBO & 19/20 & 9/20 & 0/20 & 28/60 & 0 \\
MALIBO & 12/20 & 2/20 & 0/20 & 14/60 & 0 \\
\bottomrule
\end{tabular}
}
\end{table}

No method certifies a Queue task.  These 60-task totals measure repeatability
on three fixed families, not generalization to a population of domains.

\subsection{Energy temporal audit}

\begin{table}[htbp]
\centering
\small
\caption{Postdecision Energy temporal stability.  Counts are descriptive and
do not replace the declared empirical-window certificate.}
\label{tab:temporal-audit}
\begin{tabular}{lrrrr}
\toprule
Matrix & Originally certified & Block stable & Nonoverlap stable & Joint stable \\
\midrule
Energy V2 & 317 & 303/317 & 316/317 & 303/317 \\
Functional V2 & 62 & 59/62 & 61/62 & 58/62 \\
Energy V3 & 358 & 312/358 & 358/358 & 312/358 \\
\bottomrule
\end{tabular}

\end{table}

V3 preserves all 358 original certificates under physically nonoverlapping
windows but 312 under the joint block criterion.  Shared calendars and only
five regions preclude an iid market-population interpretation.

\section{Secondary Cumulative-Variance Diagnostic}\label{app:hvd}

The project originally studied cumulative heteroscedastic variance
decomposition.  Under matched policies, the factor-HVD model represents local
and shared exposure and improves held-out variance-shape estimation.  Table
\ref{tab:hvd-diagnostic} gives the final matched diagnostic.

\begin{table}[htbp]
\centering
\small
\caption{Pooled variance versus cumulative factor-HVD.  HVD improves
calibration but not feasible recovery or verification cost.}
\label{tab:hvd-diagnostic}
\resizebox{\textwidth}{!}{\begin{tabular}{lllrrrr}
\toprule
Domain & Variance model & Feasible & Log-RMSE & Shape corr. & Mean verify \\
\midrule
FactorShock & Pooled & 20/20 & 1.158 & 0.000 & 159.2 \\
FactorShock & Cumulative factor-HVD & 20/20 & 0.219 & 0.919 & 158.4 \\
\midrule
Inventory & Pooled & 20/20 & 0.577 & 0.000 & 182.4 \\
Inventory & Cumulative factor-HVD & 20/20 & 0.349 & 0.953 & 184.0 \\
\midrule
Queue & Pooled & 19/20 & 0.501 & 0.000 & 190.4 \\
Queue & Cumulative factor-HVD & 19/20 & 0.254 & 0.995 & 197.6 \\
\bottomrule
\end{tabular}
}
\end{table}

Because both variance models recover the same feasible counts and HVD slightly
increases verification cost in Inventory and Queue, HVD is not included in the
paper's claimed optimization method.  Its formal results remain in the proof
archive as reusable mechanistic results.

\section{Machine-Checked Proof Map}\label{app:lean}

The Lean 4 project builds against mathlib without \texttt{sorry},
\texttt{admit}, or project-defined \texttt{axiom}.  Table \ref{tab:lean-map}
maps manuscript claims to their principal files.  Lean proves each displayed
implication; source-target alignment and simulator-distribution assumptions
remain scientific premises.

\begin{table}[htbp]
\centering
\small
\caption{Principal Lean files under \texttt{proof/SCOLHKG/}.}
\label{tab:lean-map}
\begin{tabular}{>{\raggedright\arraybackslash}p{0.39\textwidth}>{\raggedright\arraybackslash}p{0.52\textwidth}}
\toprule
Claim & File \\
\midrule
Finite-atlas no-free-lunch & \texttt{Real/ProposalNoFreeLunch.lean} \\
Profile-coordinate and interpolation error & \texttt{Real/ProfileCoordinateConsistency.lean} \\
Farthest-first factor-two coverage & \texttt{Real/FarthestFirstKCenter.lean} \\
Finite source-mean and rank recovery & \texttt{Measure/SourceRankRecovery.lean}, \texttt{Real/SourceRankRecovery.lean} \\
Task-law coverage calibration & \texttt{Measure/TaskAtlasCoverage.lean} \\
Aligned geometric coverage & \texttt{Real/GeometricAtlasCoverage.lean} \\
Exact all-success candidate validity & \texttt{Measure/ExactBinomialCertificate.lean} \\
Composed frontend and budget interfaces & \texttt{Real/PaperMainline.lean} \\
\bottomrule
\end{tabular}
\end{table}

\section{Reproducibility Contract}\label{app:reproducibility}

The final evidence registry contains 14 matrix contracts, 12 compact analyses,
and SHA-256 commitments for 17,360 result cells.  It records one retained
algorithmic failure and no integrity failure.  Every row identifies the task,
source archive, initial design, method, implementation commit, verifier,
source/search/verification calls, and postdecision truth reconstruction.

Tables and figures in this manuscript are generated only from compact JSON and
CSV analyses.  The renderer does not read checkpoints, pickle files, model
weights, or exported policy vectors.  Runtime artifacts are unnecessary for
auditing the displayed results.  The method specification, randomized task
law, Energy data contract, proof receipt, and rendering manifest are versioned
separately so that a theory version cannot be silently paired with another
experiment version.
